
\documentclass[a4paper,11pt]{article}
\usepackage{fontspec}
\setmainfont{Noto Sans}
\usepackage[french]{babel}
\usepackage{microtype}
\usepackage[top=9mm,bottom=8mm,left=11mm,right=11mm]{geometry}
\usepackage{xcolor}
\usepackage{tikz}
\usetikzlibrary{arrows.meta,calc,positioning}
\usepackage[most]{tcolorbox}
\pagestyle{empty}

\definecolor{Blue}{HTML}{356FA8}
\definecolor{BlueDark}{HTML}{244D78}
\definecolor{BlueLight}{HTML}{EAF2F9}
\definecolor{Border}{HTML}{B9D3E3}
\definecolor{Ink}{HTML}{172A32}
\definecolor{Grey}{HTML}{64757D}
\definecolor{Red}{HTML}{D52F35}
\definecolor{RedLight}{HTML}{FCEBED}

\tcbset{
 enhanced,boxrule=.6pt,arc=2.8mm,
 left=3.7mm,right=3.7mm,top=1.75mm,bottom=1.75mm,
 boxsep=1mm,coltext=Ink
}
\newtcolorbox{bluebox}[1][]{colback=BlueLight,colframe=Border,#1}
\newtcolorbox{warnbox}[1][]{colback=RedLight,colframe=Red,
 left=2.7mm,right=2.7mm,top=1.25mm,bottom=1.25mm,#1}

\begin{document}
\begin{tikzpicture}[remember picture,overlay]
\draw[draw=Border,line width=.9pt,rounded corners=1.7mm]
 ([xshift=1mm,yshift=-1mm]current page.north west)
 -- ([xshift=-1mm,yshift=-1mm]current page.north east)
 -- ([xshift=-1mm,yshift=1mm]current page.south east)
 -- ([xshift=1mm,yshift=1mm]current page.south west)--cycle;
\fill[Blue] ([xshift=1mm,yshift=-1mm]current page.north west)
 rectangle ([xshift=5.2mm,yshift=1mm]current page.south west);
\end{tikzpicture}

\vspace*{1mm}\hspace*{6mm}
{\color{BlueDark}\fontsize{25}{27}\selectfont\bfseries PRETERIT DE L'INDICATIU}

\vspace{0.5mm}\hspace*{6mm}
{\color{Grey}\large Le prétérit de l'indicatif}

\vspace{2.5mm}\hspace*{6mm}
\begin{minipage}{0.91\textwidth}

\begin{bluebox}
\textbf{\color{BlueDark}CONSTRUCCION}\\[1mm]
\centering
{\Large \textbf{RADICAL + TERMINASONS DEL PRETERIT}}\\[1mm]
\small
Exemple : \textbf{CANTAR}\\[1mm]
\textbf{cantèri}\quad \textbf{cantères}\quad \textbf{cantèt}\quad
\textbf{cantèrem}\quad \textbf{cantèretz}\quad \textbf{cantèron}
\end{bluebox}

\vspace{2.5mm}
\begin{center}
\begin{tikzpicture}[>=Latex]
\draw[Ink!65,line width=.8pt] (0,0)--(15.1,0);
\node[below=2mm] at (2.0,0) {\small PASSAT};
\node[below=2mm] at (7.5,0) {\small ARA};
\node[below=2mm] at (13.5,0) {\small FUTUR};
\draw[Blue,line width=2.5pt,->] (3.0,0)--(5.1,0);
\draw[BlueDark,line width=1.1pt,->] (7.5,1.0)--(7.5,.18);
\node[above=3.8mm] at (7.5,1.0) {\bfseries\color{BlueDark}ARA};
\node[above=3mm,align=center,text width=6.2cm] at (4.1,0)
{\small\bfseries eveniment acabat, dins la seguida del raconte};
\end{tikzpicture}
\end{center}

\vspace{1.5mm}
\begin{bluebox}
\textbf{\color{BlueDark}A QUÉ SERVÍS ?}\\[1mm]
\begin{itemize}\setlength{\itemsep}{0.4mm}
\item contar un \textbf{eveniment acabat} dins lo passat ;
\item far avançar la \textbf{seguida dels eveniments} d'un raconte ;
\item presentar d'accions ponctualas e successivas ;
\item s'utiliza principalament dins lo \textbf{raconte escrit e literari}.
\end{itemize}
\end{bluebox}

\vspace{2mm}
\begin{bluebox}
\textbf{\color{BlueDark}EXEMPLES}\\[1mm]
\textbf{Intrèt dins l'ostal e barrèt la pòrta.}\\
\textbf{Prenguèt lo libre, lo legiguèt puèi lo tornèt pausar.}\\
\textbf{Soudain, la pòrta se dobriguèt.}\\
\textbf{Lo rei arribèt, parlèt al pòble e partiguèt.}
\end{bluebox}

\vspace{2mm}
\begin{minipage}{0.485\textwidth}
\begin{bluebox}[height=3.25cm]
\textbf{\color{BlueDark}PRETERIT / IMPERFÈIT}\\[1mm]
Lo \textbf{preterit} presenta l'eveniment que fa avançar lo raconte.\\[1mm]
\textit{Intrèt dins la sala.}\\[1mm]
L'\textbf{imperfèit} descriu puslèu lo fons, la durada o l'abituda.
\end{bluebox}
\end{minipage}\hfill
\begin{minipage}{0.485\textwidth}
\begin{warnbox}[height=3.25cm]
\textbf{\color{Red}ATENCION}\\[1mm]
Las formas del preterit pòdon èsser \textbf{irregularas}.\\[1mm]
Cal doncas aprene las formas dels vèrbs frequents :\\
\textbf{foguèt, aguèt, faguèt, venguèt, prenguèt...}
\end{warnbox}
\end{minipage}

\vspace{2mm}
\begin{bluebox}
\textbf{\color{BlueDark}A RETÉNER}\\[1mm]
Lo preterit es lo temps del \textbf{primièr plan del raconte}: presenta d'eveniments acabats
e successius. Dins l'usatge corrent, d'autras formas del passat pòdon èsser preferidas segon lo contèxt.
\end{bluebox}

\vspace{1.5mm}\centering
{\small\color{Grey}\textbf{Repèras frequents dins los racontes :} puèi, alara, soudain, tot d'un còp, enfin, aquel jorn, quand...}

\vspace{1mm}
{\scriptsize\color{Grey}Referéncia : occitan languedocian, graphia classica, vèrb'Òc / Lo Congrès permanent de la lenga occitana.}

\end{minipage}
\end{document}
